% ------------------------------------------------
% PACKAGES
% ------------------------------------------------

\usepackage[utf8]{inputenc}
\usepackage[T1]{fontenc}

% Graphics and diagrams
\usepackage{graphicx}
\usepackage{epsfig}
\usepackage{tikz}

\usetikzlibrary{
    positioning,
    arrows.meta,
    shapes.geometric,
    shapes.misc,
    shapes.symbols,
    calc,
    fit,
    backgrounds,
    shadows
}

% Mathematics
\usepackage{amsmath}
\usepackage{amssymb}
\usepackage{latexsym}
\usepackage{amsthm}

% Header and formatting
\usepackage{fancyhdr}
\usepackage{cite}
\usepackage{setspace}
\usepackage{times}
\usepackage{microtype}
\usepackage[htt]{hyphenat}
\usepackage{xurl}
\emergencystretch=3em
\usepackage{titlesec}

\titleformat{\chapter}[display]
  {\normalfont\Large\bfseries}{\chaptertitlename\ \thechapter}{4pt}{\huge}
\titlespacing*{\chapter}{0pt}{0pt}{10pt}
\titlespacing*{\section}{0pt}{8pt plus 2pt minus 1pt}{3pt plus 1pt minus 1pt}
\titlespacing*{\subsection}{0pt}{6pt plus 1pt minus 1pt}{2pt plus 1pt minus 1pt}
\titlespacing*{\subsubsection}{0pt}{4pt plus 1pt minus 1pt}{2pt plus 1pt minus 1pt}

% Tables
\usepackage[table]{xcolor}
\definecolor{navyblue}{RGB}{15,45,95}
\usepackage{array}
\usepackage{tabularx}
\usepackage{longtable}
\usepackage{multirow}
\usepackage{booktabs}
\usepackage{blkarray}
\usepackage{bigstrut}
\usepackage{vcell}
\usepackage{adjustbox}

% Other formatting packages
\usepackage{epic}
\usepackage{float}
\usepackage{enumitem}
\setlist{itemsep=2pt,topsep=3pt,parsep=0pt,partopsep=0pt}
\usepackage{ragged2e}
\usepackage{etoolbox}
\usepackage{pdfpages}
\usepackage{imakeidx}
\usepackage{xparse}

% Code listings
\usepackage{listings}

\lstdefinelanguage{json}{
    basicstyle=\ttfamily\small,
    numbers=none,
    showstringspaces=false,
    breaklines=true,
    frame=single,
    stringstyle=\color{purple},
    keywordstyle=\color{blue},
    morekeywords={true,false,null}
}

\lstset{
    basicstyle=\ttfamily\scriptsize,
    breaklines=true,
    breakatwhitespace=false,
    columns=fullflexible,
    keepspaces=true,
    showstringspaces=false,
    frame=single,
    captionpos=b,
    numbers=none,
    tabsize=2,
    aboveskip=4pt,
    belowskip=4pt
}

% Algorithms
\usepackage{algorithm}
\usepackage{algpseudocode}

% Graph plotting
\usepackage{pgfplots}
\pgfplotsset{compat=1.18}

% Captions
\usepackage[tableposition=top]{caption}

% Hyperlinks (load near the end)
\usepackage{hyperref}

\hypersetup{
    hidelinks
}

% Index
\makeindex[columns=2]

% ------------------------------------------------
% THEOREM DEFINITIONS
% ------------------------------------------------

\newtheorem{theorem}{Theorem}[section]
\newtheorem{corollary}{Corollary}[theorem]
\newtheorem{lemma}[theorem]{Lemma}

\newfloat{procedure}{htbp}{loa}
\floatname{procedure}{Procedure}

\renewcommand{\qedsymbol}{}

\newtheorem{remark}{Remark}[section]

\algnewcommand\algorithmicinput{\textbf{Input:}}
\algnewcommand\INPUT{\item[\algorithmicinput]}

\algnewcommand\algorithmicoutput{\textbf{Output:}}
\algnewcommand\OUTPUT{\item[\algorithmicoutput]}

\algnewcommand{\algorithmicand}{\textbf{ and }}
\algnewcommand{\algorithmicor}{\textbf{ or }}
\algnewcommand{\OR}{\algorithmicor}
\algnewcommand{\AND}{\algorithmicand}

\newtheorem{thm}{Theorem}[section]
\newtheorem{lem}[thm]{Lemma}
\newtheorem{prop}[thm]{Proposition}
\newtheorem{cor}{Corollary}
\newtheorem{conj}{Conjecture}[section]

\theoremstyle{definition}
\newtheorem{defn}{Definition}[section]
\newtheorem{exmp}{Example}[section]
\newtheorem{rem}{Remark}

% ------------------------------------------------
% SECTION AND TOC DEPTH
% ------------------------------------------------

\setcounter{secnumdepth}{3}
\setcounter{tocdepth}{1}

% Float parameters to allow figures and text to share pages
\renewcommand{\topfraction}{0.92}
\renewcommand{\bottomfraction}{0.85}
\setcounter{topnumber}{4}
\setcounter{bottomnumber}{3}
\setcounter{totalnumber}{6}
\renewcommand{\textfraction}{0.05}
\renewcommand{\floatpagefraction}{0.75}

\setlength{\textfloatsep}{8pt plus 2pt minus 2pt}
\setlength{\intextsep}{6pt plus 2pt minus 2pt}
\setlength{\floatsep}{6pt plus 2pt minus 2pt}
\captionsetup{skip=4pt}

% ------------------------------------------------
% BIBLIOGRAPHY NAME
% ------------------------------------------------

\renewcommand{\bibname}{References}

% ------------------------------------------------
% LINE SPACING & PARAGRAPH FORMATTING
% ------------------------------------------------

\renewcommand{\baselinestretch}{1.08}
\setlength{\parskip}{2pt}
\setlength{\parindent}{1.2em}

% ------------------------------------------------
% PAGE HEADERS AND FOOTERS
% ------------------------------------------------

\pagestyle{fancy}

\renewcommand{\chaptermark}[1]{
    \markright{\chaptername\ \thechapter.\ #1}
}

\renewcommand{\sectionmark}[1]{
    \markright{\thesection.\ #1}
}

\lhead[\fancyplain{}{}]
      {\fancyplain{}{\bfseries\rightmark}}

\chead[\fancyplain{}{}]
      {\fancyplain{}{}}

\rhead[\fancyplain{}{}]
      {\fancyplain{}{\bfseries\thepage}}

\lfoot[\fancyplain{}{}]
      {\fancyplain{}{}}

\cfoot[\fancyplain{}{}]
      {\fancyplain{}{}}

\rfoot[\fancyplain{}{}]
      {\fancyplain{}{\bfseries\today}}

% ------------------------------------------------
% FOOTNOTE SPACING
% ------------------------------------------------

\addtolength{\skip\footins}{5mm}

% ------------------------------------------------
% EQUATION NUMBERING
% ------------------------------------------------

\newcounter{ind}

\def\eqlabon{
    \setcounter{ind}{\value{equation}}
    \addtocounter{ind}{1}
    \setcounter{equation}{0}
    \renewcommand{\theequation}{
        \arabic{chapter}.\arabic{section}.
        \arabic{ind}\alph{equation}
    }
}

\def\eqlaboff{
    \renewcommand{\theequation}{
        \arabic{chapter}.\arabic{section}.\arabic{equation}
    }
    \setcounter{equation}{\value{ind}}
}

% ------------------------------------------------
% MATHEMATICAL SHORTCUTS
% ------------------------------------------------

\def\bv#1{\mbox{\boldmath$#1$}}

\def\mint{\textstyle\mints\displaystyle}
\def\mints{\int\!\!\!\!\!\!{\rm-}\ }

\def\msum{\textstyle\msums\displaystyle}
\def\msums{\sum\!\!\!\!\!\!\!{\rm-}\ }

\def\poincare{Poincar\'e }
\def\holder{H\"older }